% C-附录C-常见机器人动力学.tex
\documentclass[UTF8]{ctexart}
\usepackage{amsmath,amssymb,amsfonts}
\usepackage{geometry}
\geometry{a4paper,margin=2.5cm}
\usepackage{graphicx}
\usepackage{float}

\begin{document}

\section*{附录C}
\addcontentsline{toc}{section}{附录C 常见机器人机械臂的动力学}

\subsection*{常见机器人机械臂的动力学}

在本附录中，我们给出一些常见机器人机械臂的动力学。假设机器人动力学由下式给出：

\begin{equation}
M(q)\ddot{q} + N(q, \dot{q}) = \tau
\tag{C.1.1}
\end{equation}

其中矩阵 $M(q)$ 是对称且正定的，其元素为 $m_{ij}(q)$，即

\begin{equation}
M(q) = [m_{ij}(q)], \quad i,j = 1, \ldots, n
\end{equation}

而 $N(q, \dot{q})$ 是一个 $n \times 1$ 向量，其元素为 $n_i$，即

\begin{equation}
N(q, \dot{q}) = [n_i], \quad i = 1, \ldots, n
\end{equation}

特别要注意，在 $n_i$ 的表达式中通过重力常数 $g = 9.8$ 米/秒$^2$ 来识别重力项。我们还将采用以下符号表示：

\begin{itemize}
\item 连杆 $i$ 的长度为 $L_i$，单位为米
\item 连杆 $i$ 的质量为 $m_i$，单位为千克
\item 连杆 $i$ 关于轴 $u$ 的质量惯性矩为 $I_{uui}$，单位为 kg-m$^2$
\item $S_i = \sin q_i$，$C_i = \cos q_i$
\item $S_{ij} = \sin(q_i + q_j)$，$C_{ij} = \cos(q_i + q_j)$
\item $S_{ijk} = \sin(q_i + q_j + q_k)$，$C_{ijk} = \cos(q_i + q_j + q_k)$
\item $SS_i = \sin^2 q_i$，$CC_i = \cos^2 q_i$，$CS_i = \cos q_i \sin q_i$
\item $SS_{ij} = \sin^2(q_i + q_j)$
\end{itemize}

\subsection*{C.1 SCARA 机械臂}

我们考虑的第一个机器人是通用 SCARA 构型机器人，如图 C.1.1 所示。这些方程适用于 AdeptOne 和 AdeptTwo 机器人。动力学包括前四个自由度，以符号形式给出如下：

\begin{align}
m_{11} &= m_1L_1^2 + I_{zz1} + m_2(L_1^2 + L_2^2 + 2L_1L_2C_2) + I_{zz2} + m_3(L_1^2 + L_2^2 + 2L_1L_2C_2) \\
m_{12} &= m_2(L_2^2 + L_1L_2C_2) + I_{zz2} + m_3(L_2^2 + L_1L_2C_2) \\
m_{13} &= 0 \\
m_{14} &= 0 \\
m_{22} &= m_2L_2^2 + I_{zz2} + m_3L_2^2 \\
m_{23} &= 0 \\
m_{24} &= 0 \\
m_{33} &= m_3 \\
m_{34} &= 0 \\
m_{44} &= I_{zz4}
\end{align}

\begin{align}
n_1 &= -m_2L_1L_2S_2(2\dot{q}_1\dot{q}_2 + \dot{q}_2^2) - m_3L_1L_2S_2(2\dot{q}_1\dot{q}_2 + \dot{q}_2^2) \\
n_2 &= m_2L_1L_2S_2\dot{q}_1^2 + m_3L_1L_2S_2\dot{q}_1^2 \\
n_3 &= -m_3g \\
n_4 &= 0
\end{align}

\begin{figure}[H]
\centering
\fbox{\parbox{0.8\textwidth}{\centering 图 C.1.1：SCARA 机械臂}}
\caption{SCARA 机械臂示意图}
\end{figure}

\subsection*{C.2 Stanford 机械臂}

如图 C.2.1 所示的 Stanford 机械臂具有以下动力学方程 [Bejczy 1974], [Paul 1981]：

\begin{align}
m_{11} &= I_{xx2}S_2^2 + I_{yy2}C_2^2 + m_2d_2^2C_2^2 + m_3q_3^2S_2^2 + I_{xx3}S_2^2 + I_{yy3}C_2^2 \\
m_{12} &= -m_3q_3C_2 \\
m_{13} &= -m_3d_2S_2 \\
m_{22} &= m_3q_3^2 + I_{zz2} + I_{zz3} \\
m_{23} &= 0 \\
m_{33} &= m_3
\end{align}

\begin{align}
n_1 &= (I_{yy2} - I_{xx2} + m_3q_3^2 + I_{yy3} - I_{xx3})CS_2\dot{q}_1\dot{q}_2 + m_3q_3^2S_2C_2\dot{q}_1\dot{q}_3 \\
&\quad + (I_{yy3} - I_{xx3} + m_3q_3^2)CS_2\dot{q}_1\dot{q}_2 \\
n_2 &= -(I_{yy2} - I_{xx2} + m_3q_3^2 + I_{yy3} - I_{xx3})CS_2\dot{q}_1^2/2 - m_3q_3^2S_2C_2\dot{q}_1^2/2 \\
&\quad + m_3q_3S_2\dot{q}_2\dot{q}_3 - m_3q_3C_2^2\dot{q}_1^2 - m_3q_3\dot{q}_2^2/2 + m_3gC_2 \\
n_3 &= -m_3q_3S_2^2\dot{q}_1^2 - m_3d_2C_2\dot{q}_2^2 + m_3q_3\dot{q}_2^2 - m_3gS_2
\end{align}

\begin{figure}[H]
\centering
\fbox{\parbox{0.8\textwidth}{\centering 图 C.2.1：Stanford 机械臂}}
\caption{Stanford 机械臂示意图}
\end{figure}

\subsection*{C.3 PUMA 560 机械臂}

PUMA 560 如图 C.3.1 所示。对于这种特定的结构，可以进行许多简化以获得以下动力学方程，这些方程出现在 [Armstrong et al. 1986] 中。

\begin{align}
m_{11} &= I_{xx1} + I_{xx2}S_2^2 + I_{yy2}C_2^2 + m_2d_2^2C_2^2 + m_3(L_2S_2 + d_3C_2)^2 \\
&\quad + I_{xx3}S_{23}^2 + I_{yy3}C_{23}^2 + m_3d_3^2S_{23}^2 + m_4(L_2S_2 + L_3S_{23})^2 \\
m_{12} &= m_3d_2d_3S_3 + m_4d_2L_3S_3 \\
m_{13} &= -m_3d_3^2 - m_3L_2d_3C_3 - m_4L_2L_3C_3 - I_{zz3} \\
m_{22} &= m_2d_2^2 + m_3d_3^2 + m_3L_2^2 + m_4L_2^2 + m_4L_3^2 + 2m_3L_2d_3C_3 + 2m_4L_2L_3C_3 + I_{zz2} + I_{zz3} \\
m_{23} &= m_3d_3^2 + m_4L_3^2 + m_3L_2d_3C_3 + m_4L_2L_3C_3 + I_{zz3} \\
m_{33} &= m_3d_3^2 + m_4L_3^2 + I_{zz3}
\end{align}

\begin{align}
n_1 &= [(I_{yy2} - I_{xx2})C_2S_2 + m_2d_2^2C_2S_2 + m_3(L_2S_2 + d_3C_2)(L_2C_2 - d_3S_2) \\
&\quad + (I_{yy3} - I_{xx3})C_{23}S_{23} + m_3d_3^2C_{23}S_{23} + m_4(L_2S_2 + L_3S_{23})(L_2C_2 + L_3C_{23})](2\dot{q}_1\dot{q}_2) \\
&\quad + [(I_{yy3} - I_{xx3})C_{23}S_{23} + m_3d_3^2C_{23}S_{23} - m_3(L_2S_2 + d_3C_2)d_3S_{23} \\
&\quad + m_4(L_2S_2 + L_3S_{23})L_3C_{23}](2\dot{q}_1\dot{q}_3) \\
n_2 &= [(I_{xx2} - I_{yy2})C_2S_2 - m_2d_2^2C_2S_2 - m_3(L_2S_2 + d_3C_2)(L_2C_2 - d_3S_2) \\
&\quad + (I_{xx3} - I_{yy3})C_{23}S_{23} - m_3d_3^2C_{23}S_{23} - m_4(L_2S_2 + L_3S_{23})(L_2C_2 + L_3C_{23})]\dot{q}_1^2 \\
&\quad + [-m_3L_2d_3S_3 - m_4L_2L_3S_3](2\dot{q}_2\dot{q}_3 + \dot{q}_3^2) \\
&\quad + [m_3d_3C_{23} + m_4L_3C_{23}]g \\
n_3 &= [(I_{xx3} - I_{yy3})C_{23}S_{23} - m_3d_3^2C_{23}S_{23} + m_3(L_2S_2 + d_3C_2)d_3S_{23} \\
&\quad - m_4(L_2S_2 + L_3S_{23})L_3C_{23}]\dot{q}_1^2 \\
&\quad + [m_3L_2d_3S_3 + m_4L_2L_3S_3]\dot{q}_2^2 \\
&\quad + [m_3d_3S_{23} + m_4L_3S_{23}]g
\end{align}

\begin{figure}[H]
\centering
\fbox{\parbox{0.8\textwidth}{\centering 图 C.3.1：PUMA 560 机械臂}}
\caption{PUMA 560 机械臂示意图}
\end{figure}

\vspace{1em}

\subsection*{参考文献}

\noindent [Armstrong et al. 1986] Armstrong, B., O. Khatib, and J. Burdick, ``The explicit dynamic model and inertial parameters of the PUMA 560 arm,'' \textit{Proc. 1986 IEEE Conf. Robot. Autom.}, pp. 510--518, San Francisco, Apr. 7--10, 1986.

\vspace{0.5em}

\noindent [Bejczy 1974] Bejczy, A.K., ``Robot arm dynamics and control,'' NASA-JPL Technical Memorandum 33--669, 1974.

\vspace{0.5em}

\noindent [Paul 1981] Paul, R.P., \textit{Robot Manipulators: Mathematics, Programming and Control}. Cambridge, MA: MIT Press, 1981.

\end{document}
